\section{Seeds and one algorithm for all sizes}
\label{sec:seeds}

This section turns the relative scheme of Theorem~\ref{thm:main} into a
rational seed (Section~\ref{sec:seeds-seed}), bounds its coefficients and
the cost of generating them (Section~\ref{sec:seeds-coefficients}), and
gives one algorithm for all sizes (Section~\ref{sec:seeds-uniform}).

\paragraph{Setting.}
Throughout, \(\eps\in(0,1]\) is rational, \(\Lambda=\log(4096/\eps)\),
and \(\mathcal S_\eps\) is the relative scheme \(r_*\MM_A\restr m_*\MM_B\)
of Theorem~\ref{thm:main}, with the parameters \(D,R,L,H\) and \(N=k\)
defined there. Its chain has \(J=2(R-1)+L\) stages, from the source pair
\((r_1,b_1)\) to the terminal pair \((1,1)\), with arities
\(2\le d_i<24D\). At stage \(i\), \(\mathcal S_\eps\) uses the first type,
in the lexicographic order of \((\nK,\nO)\), with the largest sector count
\(M_i\); so \((d_i+2)^N/(N+1)^2\le M_i<(d_i+2)^N\)
(Lemma~\ref{lem:types}). The counts are those
of~\eqref{eq:stage-chain-counts} for \(t=1\), and \(s\) and \(E_i\)
are as in Proposition~\ref{prop:chain-counts}; the three maps are
rational, and \(9/4<p<9/4+\eps/2\).

\paragraph{The pooled recursion.}
We use the pooled recursion of \citep[Lemma~2.4]{FinCon26} in its
notation, except that we write \(\bar n_\ell\) for the side at depth
\(\ell\), because \(n_0\) denotes the side of a seed here. Thus
\(\rho=A/B\), \(n_\infty=AB/(B-A)\), and \(n_*=\lceil n_\infty\rceil+1\)
is the leaf threshold. At side \(n\), the recursion starts from
\(\nu_0=1\) product of side \(\bar n_0=n\) and sets
\(\nu_{\ell+1}=r_*\lceil\nu_\ell/m_*\rceil\) and
\(\bar n_{\ell+1}=A\lceil\bar n_\ell/B\rceil\); the last depth
\(\ell_{\max}\) is the first \(\ell\) with \(\bar n_\ell\le n_*\), and its
products are computed classically. Step~3 of the proof of
\citep[Lemma~2.4]{FinCon26} shows that, for every \(n\ge1\) and
\(0\le\ell\le\ell_{\max}\),
\begin{equation}\label{eq:seeds-pooled}
 \nu_\ell\le\Bigl(1+\frac{r_*m_*}{r_*-m_*}\Bigr)\Bigl(\frac{r_*}{m_*}\Bigr)^\ell,
 \qquad
 \rho^\ell n\le\bar n_\ell<n_\infty+\rho^\ell n,
 \qquad
 \rho<\rho^{\ell_{\max}}n\le n_*.
\end{equation}

\subsection{Resources}
\label{sec:seeds-resources}

\begin{lemma}[Resources]\label{lem:resources}
For every rational \(\eps\in(0,1]\), the scheme \(\mathcal S_\eps\)
satisfies:
\begin{enumerate}
\item[(a)] \(N=k\ge4608\), \(J\le240\,\eps^{-3/2}\),
 \(N\le2305\,\eps^{-3/2}\Lambda\), and
 \(2NkJ\le\log_2A\le2NkJ\log_2(24D)\); in particular
 \(\log A\le2{,}550{,}252{,}000\,\eps^{-9/2}\Lambda^3\);
\item[(b)] \(m_*=B^3\) and \(\log(B/A)\ge124Nk\); in particular
 \(2A\le B\le A^2\);
\item[(c)] \(4\le(B/A)^2\le r_*/m_*\le(B/A)^{11/4}\) and
 \(r_*<A^{35/4}\);
\item[(d)] \(A<n_*\le2A+1\);
\item[(e)] \(A^{12}\le\Xi\) and \(\Xi+1\le A^{27}\), where
 \(\Xi=r_*(1+r_*m_*)\,n_*^3\).
\end{enumerate}
\end{lemma}

\begin{proof}
(a) Steps~0 and~5 of the proof of Theorem~\ref{thm:main} give
\(N=k\ge4608\), \(J\le240\,\eps^{-3/2}\),
\(N\le2305\,\eps^{-3/2}\Lambda\) and the bound on \(\log A\). Since
\(2\le d_i<24D\), the number \(\log_2A=2Nk\sum_i\log_2d_i\) lies between
\(2NkJ\) and \(2NkJ\log_2(24D)\).

(b) Theorem~\ref{thm:main} gives \(m_*=B^3\) and \(B\le A^2\). Let
\(G=\sum_i\log((d_i+2)/d_i)\) be the gain of the chain. Step~1 of the
proof of Theorem~\ref{thm:main} gives
\begin{equation}\label{eq:seeds-gain}
 \log(B/A)\ge\frac{497}{512}\cdot2NkG,\qquad
 G\ge\frac{64}\eps,\qquad J\le\frac{5D}4\,G,
\end{equation}
so \(\log(B/A)\ge2Nk\cdot\frac{64}\eps\cdot\frac{497}{512}\ge124Nk\), and
in particular \(B\ge2A\).

(c) By Step~1 of the proof of \citep[Lemma~2.4]{FinCon26}, \(p\ge2\), so
\(r_*/m_*=(B/A)^p\ge(B/A)^2\ge4\). Also \(p<9/4+\eps/2\le11/4\), and
\(r_*<A^{35/4}\) by Theorem~\ref{thm:main}.

(d) Since \(B\ge2A\), we have \(A<AB/(B-A)=A/(1-A/B)\le2A\), so
\(A+2\le n_*\le2A+1\).

(e) By (b)--(d), \(r_*\ge4m_*\), \(m_*=B^3>A^3\) and \(n_*>A\), so
\(\Xi\ge r_*^2m_*n_*^3>A^{12}\). Also \(1+r_*m_*\le2r_*m_*\) and
\(n_*\le3A\), so \(\Xi\le54\,r_*^2m_*A^3\le54A^{35/2+6+3}\). Hence
\(\Xi+1\le55A^{53/2}\le A^{27}\), because \(A\ge2^{2NkJ}>55^2\).
\end{proof}

\subsection{An explicit seed}
\label{sec:seeds-seed}

The next theorem is the analogue of \citep[Theorem~9.2]{FinCon26} for
\(\mathcal S_\eps\), and its proof follows the proof given there.

\begin{theorem}[An explicit seed]\label{thm:seed}
Let \(\eps\in(0,1]\) be rational, let \(n_*\) and \(\Xi\) be as in
Lemma~\ref{lem:resources}, and put
\[
 n_0=(\Xi+1)^{\lceil2/\eps\rceil}.
\]
For every \(n\ge n_0\), unrolling the pooled recursion of
\(\mathcal S_\eps\) at side \(n\), with the zero padding of
\citep[Lemma~2.4]{FinCon26} and classical multiplication at the last
depth, gives a decomposition of \(\MM_n\) into at most
\(\Xi n^p<n^{9/4+\eps}\) rank-one tensors with rational coefficients. In
particular, \(n=n_0\) gives a rational seed of side \(n_0\) with fewer
than \(n_0^{9/4+\eps}\) products, and
\[
 \frac{24}\eps\log A\;\le\;\log n_0\;\le\;
 27\Bigl\lceil\frac2\eps\Bigr\rceil\log A
 \;<\;2.1\cdot10^{11}\,\eps^{-11/2}\log^3\frac{4096}\eps .
\]
\end{theorem}

\begin{proof}
Run the pooled recursion at side \(n\). By~\eqref{eq:seeds-pooled} and
\(r_*-m_*\ge1\), we have
\(\nu_{\ell_{\max}}\le(1+r_*m_*)(r_*/m_*)^{\ell_{\max}}\). Since
\(r_*/m_*=\rho^{-p}\) and \(\rho^{\ell_{\max}}n>\rho\),
\[
 \Bigl(\frac{r_*}{m_*}\Bigr)^{\ell_{\max}}=(\rho^{\ell_{\max}})^{-p}
 <\Bigl(\frac n\rho\Bigr)^p=\frac{r_*}{m_*}\,n^p,
\]
and therefore \(\nu_{\ell_{\max}}<r_*(1+r_*m_*)\,n^p\), as \(m_*\ge1\).
Every product at the last depth has side
\(\bar n_{\ell_{\max}}\le n_*\), and classical multiplication computes it
with at most \(n_*^3\) scalar products. Each of them multiplies an entry
of a left input by an entry of a right input, which by Step~5 of the proof
of \citep[Lemma~2.4]{FinCon26} are linear forms in the entries of the two
original \(n\times n\) matrices. All later operations (the output maps and
the removal of padding) are fixed linear maps, so every entry of the
computed product is a fixed linear combination of these scalar products.
As in the proof of \citep[Corollary~8.7]{FinCon26}, the scalar products
therefore give a decomposition of \(\MM_n\) into at most
\(\nu_{\ell_{\max}}n_*^3<\Xi n^p\) rank-one tensors, one for each scalar
product. Its coefficients are rational because the maps of
\(\mathcal S_\eps\) are rational. Finally, \(n\ge n_0\) gives
\(n^{\eps/2}\ge(\Xi+1)^{\lceil2/\eps\rceil\eps/2}\ge\Xi+1>\Xi\), and
\(p<9/4+\eps/2\) gives \(\Xi n^p<n^{9/4+\eps}\).

For the size, \(\log n_0=\lceil2/\eps\rceil\log(\Xi+1)\), and
Lemma~\ref{lem:resources}(e) gives
\((2/\eps)\cdot12\log A\le\log n_0\le27\lceil2/\eps\rceil\log A\). Since
\(\eps\le1\), we have \(\lceil2/\eps\rceil\le3/\eps\), and
Lemma~\ref{lem:resources}(a) gives
\(\log n_0\le81\cdot2{,}550{,}252{,}000\,\eps^{-11/2}\Lambda^3
<2.1\cdot10^{11}\,\eps^{-11/2}\Lambda^3\).
\end{proof}

Block recursion with the seed is covered by \citep[Lemma~9.4 and
Propositions~9.5--9.6]{FinCon26}.

\begin{remark}[The factor \(1/\eps\)]\label{rem:seed-factor}
The factor of order \(1/\eps\) between \(\log n_0\) and \(\log A\) is
necessary for the procedure of Theorem~\ref{thm:seed} when
\(\eps\le3/4\), as in \citep[Remark~8.8]{FinCon26}: if the procedure,
run at side \(n\), executes fewer than \(n^{9/4+\eps}\) scalar products
at its last depth, counting those in batch slots filled with zero
matrices, then \(\log n>(\frac{13}4\log A-\log27)/\eps\). Indeed,
\(\ell_{\max}=0\) would give \(n^3\ge n^{9/4+\eps}\) products, so
\(\ell_{\max}\ge1\). Then \(\nu_1=r_*\) and
\(\nu_{\ell+1}\ge(r_*/m_*)\nu_\ell\), so
\(\nu_{\ell_{\max}}\ge m_*\rho^{-p\ell_{\max}}\ge m_*(n/n_*)^p\)
by~\eqref{eq:seeds-pooled}. Every product at the last depth has side
\(A\lceil\bar n_{\ell_{\max}-1}/B\rceil\ge A\), hence needs at least
\(A^3\) scalar products. With \(m_*=B^3\ge A^3\), \(n_*\le3A\) and
\(p\le3\), the procedure executes at least \(n^pA^{6-p}/27\) products,
and \(9/4<p<21/8\) (Lemma~\ref{lem:profile-potential} and
Theorem~\ref{thm:main}) gives
\(\eps\log n>(9/4+\eps-p)\log n>(6-p)\log A-\log27
>\frac{13}4\log A-\log27\). This does not bound the rank after rank-one
terms with a zero factor are deleted.
\end{remark}

\subsection{Coefficients and generation}
\label{sec:seeds-coefficients}

Appendix~\ref{sec:coefficients} writes each of the three maps of
\(\mathcal S_\eps\) as a product of explicit matrices and describes a
simple generator that multiplies them out. We count bit operations of
schoolbook integer arithmetic; an addition, subtraction, multiplication,
comparison or division with remainder of integers with at most
\(b\ge1\) bits costs at most \(16b^2\) of them.

\begin{proposition}[Coefficients and generation]\label{prop:coefficients}
Let \(\eps\in(0,1]\) be rational, and put
\[
 \mathcal D=((s-1)!)^{2Nk}\prod_{i=1}^J
 \bigl((M_i-1)!\,(2M_i-2)!\bigr)^{2k}\,(E_i!)^{2k(M_i-1)},\qquad
 \beta_\eps=(24D)^{3N}.
\]
\begin{enumerate}
\item[(a)] Every coefficient of the three maps of \(\mathcal S_\eps\) is
 \(\mathfrak c/\mathcal D\) with an integer \(\mathfrak c\),
 \(|\mathfrak c|<2^{\beta_\eps}\), and \(\mathcal D<2^{\beta_\eps}\).
\item[(b)] \(\beta_\eps\le A^{1/k}\), and
 \(\log_2\beta_\eps<10^4\,\eps^{-3/2}\Lambda^2\).
\item[(c)] Let \(\eps=u/v\) with positive integers \(u,v\le A\); this
 holds whenever \(\eps=1/v\), since \(v<N<A\). The generator of
 Appendix~\ref{sec:coefficients-generator} computes \(\mathcal D\) and
 the integers \(\mathfrak c\) of all three maps from \(u\) and \(v\) with
 at most \(A^{46}\) bit operations, and every integer that it forms has
 absolute value below \(2^{\beta_\eps}\).
\end{enumerate}
\end{proposition}

\begin{proof}
Parts (a) and (c) are proved in
Appendices~\ref{sec:coefficients-bounds}
and~\ref{sec:coefficients-generator}. For (b), \(\log_2A^{1/k}
=2N\sum_i\log_2d_i\ge2NJ\), and \(2J\ge2L\ge128D\ge3\log_2(24D)\), so
\(\beta_\eps\le A^{1/k}\). Also \(\log_2\beta_\eps=3N\log_2(24D)
\le3\cdot2305\,\eps^{-3/2}\Lambda\cdot\Lambda/\log2<10^4\,\eps^{-3/2}\Lambda^2\),
by Lemma~\ref{lem:resources}(a) and \(\log(24D)\le\Lambda\) (Step~5 of
the proof of Theorem~\ref{thm:main}).
\end{proof}

The power \(46\) in (c) comes from dense products of matrices of side
at most \(r_*A^2B^2\); the generator exploits no sparsity.

\begin{remark}[Precision of one batch]\label{rem:precision}
Apply \(\mathcal S_\eps\) once to \(m_*\) pairs of \(B\times B\)
matrices \(X_j,Y_j\), and multiply the \(r_*\) pairs of \(A\times A\)
matrices that it forms classically. Evaluate every linear combination by
recursive summation in binary floating point with round-to-nearest and
unit roundoff \(\mathsf u\le2^{-4\beta_\eps}\), without underflow or
overflow. Then every entry of every computed product \(X_jY_j\) has error
at most \(2^{4\beta_\eps}\mathsf u\max_j\|X_j\|_{\max}\max_j\|Y_j\|_{\max}\).
The proof follows that of \citep[Proposition~9.6]{FinCon26}, for one
level with classical leaves of side~\(A\). A bilinear monomial carries at
most \(\chi=2m_*B^2+r_*A^2+A+3\) rounding factors, and
\(\chi\le4A^{43/4}\) by Lemma~\ref{lem:resources}(b),(c). An output entry
combines at most \(r_*A^2\) entries of the \(A\times A\) products, each
a sum of \(A\) products of two linear combinations, and the leg matrices
have entries and absolute row sums below \(2^{\beta_\eps}\)
(Appendix~\ref{sec:coefficients-bounds}). So the exact monomials of one
entry have absolute sum at most \(\Upsilon=r_*A^32^{3\beta_\eps}\) times
\(\max_j\|X_j\|_{\max}\max_j\|Y_j\|_{\max}\). As
\(2\chi\Upsilon\le8A^{45/2}2^{3\beta_\eps}<2^{4\beta_\eps}\), we have
\(\chi\mathsf u<1/2\), and the error is
at most \(\Upsilon\chi\mathsf u/(1-\chi\mathsf u)\le2^{4\beta_\eps}\mathsf u\)
times that product. We do not analyze the precision or the integer bit
complexity of the pooled recursion, in which the inputs of a batch are
themselves computed.
\end{remark}

\subsection{One algorithm for all sizes}
\label{sec:seeds-uniform}

We first make the constants of the pooled recursion explicit.

\begin{lemma}[Pooled recursion with explicit constants]
\label{lem:pooled-constants}
Let \(r_*\MM_A\restr m_*\MM_B\) be a relative scheme with pooled exponent
\(p\le3\), \(B\ge2A\) and \(m_*\le B^3\). Store its three maps as dense
matrices over a field, and run the pooled recursion of
\citep[Lemma~2.4]{FinCon26} with classical multiplication at the last
depth. For every \(n\ge1\), the number of additions, subtractions and
multiplications in the field, including multiplications by coefficients
of the scheme, is at most
\[
 C_\nu n^p\Bigl[2\frac{r_*}{m_*}n_*^3+12m_*B^2\Bigl(\frac43\frac{r_*}{m_*}
 n_\infty^2+\ell_{\max}\Bigl(\frac BA\Bigr)^{p-2}\Bigr)\Bigr]
 \le(304+28\log_2n)\,B^{11}n^p,
\]
where \(C_\nu=1+r_*m_*/(r_*-m_*)\) and \(\ell_{\max}\) is the last depth at
side \(n\).
\end{lemma}

\begin{proof}
Let \(\ell<\ell_{\max}\) and \(c_\ell=\lceil\bar n_\ell/B\rceil\). One
batch at depth \(\ell\) is an instance of \(m_*\MM_B\) whose entries are
blocks of side \(c_\ell\). Each of the \(r_*A^2\) left blocks of the
created products is a combination of at most \(m_*B^2\) blocks, and so is
each right block; each of the \(m_*B^2\) output blocks is a combination of
at most \(r_*A^2\) blocks of the products. So a batch costs fewer than
\(6r_*m_*A^2B^2c_\ell^2\) operations. Since
\(\nu_{\ell+1}=r_*\lceil\nu_\ell/m_*\rceil\) and
\(\bar n_{\ell+1}=Ac_\ell\), depth \(\ell\) costs at most
\(\lceil\nu_\ell/m_*\rceil\cdot6r_*m_*A^2B^2c_\ell^2
=6m_*B^2\nu_{\ell+1}\bar n_{\ell+1}^2\). The last depth costs at most
\(2\nu_{\ell_{\max}}n_*^3\).

By~\eqref{eq:seeds-pooled}, \(\nu_\ell\le C_\nu(r_*/m_*)^\ell\) and
\(\bar n_\ell^2\le2n_\infty^2+2\rho^{2\ell}n^2\). Since
\(r_*/m_*=(B/A)^p\ge4\) and \(\rho^{\ell_{\max}}n>\rho\),
\[
 \sum_{\ell=1}^{\ell_{\max}}\Bigl(\frac{r_*}{m_*}\Bigr)^\ell
 \le\frac43\Bigl(\frac{r_*}{m_*}\Bigr)^{\ell_{\max}}
 <\frac43\,\frac{r_*}{m_*}\,n^p .
\]
The terms \(\rho^{(2-p)\ell}n^2=(r_*/m_*)^\ell\rho^{2\ell}n^2\) do not
decrease with \(\ell\), and the last one is
\((\rho^{\ell_{\max}}n)^{2-p}n^p\le\rho^{2-p}n^p\), so their sum over
\(1\le\ell\le\ell_{\max}\) is at most \(\ell_{\max}(B/A)^{p-2}n^p\).
Adding the leaves, which cost at most
\(2C_\nu(r_*/m_*)^{\ell_{\max}}n_*^3<2C_\nu(r_*/m_*)n^pn_*^3\), gives the
first bound.

For the second, \(r_*\ge4m_*\) gives
\(C_\nu\le1+\frac43m_*\le\frac73B^3\). Also \(r_*/m_*\le(B/A)^3\),
\(n_*\le2A+1\le3A\), \(n_\infty\le2A\) and \((B/A)^{p-2}\le B/A\). Since
\(\rho\le1/2\) and \(\rho^{\ell_{\max}-1}n>1\), we have
\(\ell_{\max}\le\log_2n+1\). The three terms are then at most
\(126B^6\), \(\frac{448}3B^{11}/A\) and \(28\ell_{\max}B^9/A\), and their sum
is at most \((126+150+28(\log_2n+1))B^{11}\).
\end{proof}

The lemma counts field operations, as \citep[Corollary~9.7]{FinCon26}
does. Padding, cutting and reassembling at depth \(\ell\) move at most
\(3\lceil\nu_\ell/m_*\rceil m_*(Bc_\ell)^2\) entries, at most half the
operations counted for that depth.

For an integer \(v\ge1\), let \(D_v\), \(L_v\), \(H_v\), \(N_v\) and
\(J_v\) be \(D\), \(L\), \(H\), \(N\) and \(J\) for \(\eps=1/v\), and put
\[
 \mathsf a(v)=2N_v^2J_v\cdot\bigl(\text{number of binary digits of }24D_v\bigr).
\]
By Lemma~\ref{lem:resources}(a), \(\log_2A\le\mathsf a(v)\) for the scheme
\(\mathcal S_{1/v}\). Each of \(D_v\), \(L_v\), \(H_v\), \(N_v\) and \(J_v\)
is nondecreasing in \(v\), so \(\mathsf a\) is nondecreasing, and
\(\mathsf a(1)=2\cdot4608^2\cdot158\cdot6=40{,}259{,}026{,}944\).

\begin{definition}[Algorithm \(\mathcal U\)]\label{def:algorithm-uniform}
The input is an integer \(n\ge1\) and two \(n\times n\) matrices over a
field of characteristic zero.
\begin{enumerate}
\item Compute \(\kappa=\lfloor\log_2n\rfloor\). If
 \(3000\,\mathsf a(1)>\kappa\), multiply the matrices classically and
 stop. Otherwise compute
 \(v(n)=\max\{v\ge1:3000\,v\,\mathsf a(v)\le\kappa\}\) by binary search
 over \(1\le v\le\kappa\).
\item Run the generator of Appendix~\ref{sec:coefficients-generator} for
 \(\eps=1/v(n)\). It outputs \(\mathcal D\) and the integer numerators of
 the three maps of \(\mathcal S_{1/v(n)}\).
\item Convert \(\mathcal D\) and every numerator to the field by Horner's
 rule in base~\(2\), invert \(\mathcal D\), and multiply every numerator
 by \(\mathcal D^{-1}\).
\item Multiply the matrices by the pooled recursion of
 \(\mathcal S_{1/v(n)}\) with classical leaves, as in
 Lemma~\ref{lem:pooled-constants}.
\end{enumerate}
\end{definition}

\begin{theorem}[One algorithm for all sizes]\label{thm:uniform}
Count steps as in \citep[Corollary~9.7]{FinCon26}: a step is one step of a
fixed multitape Turing machine during the preprocessing (Steps~1 and~2 of
Algorithm~\(\mathcal U\), and the reading of the generator's output in
Step~3), or one addition, subtraction, multiplication or division in the
field on matrix entries or scheme coefficients. There is a constant
\(c_0\), depending only on the fixed machine, with the following
properties. Algorithm~\(\mathcal U\) multiplies two \(n\times n\) matrices
over a field of characteristic zero. If
\(\lfloor\log_2n\rfloor\ge3000\,\mathsf a(1)\approx1.21\cdot10^{14}\), it
uses at most
\[
 (c_0+28\log_2n)\,n^{9/4+0.51/v(n)}
\]
steps; otherwise it uses at most \(2n^3+c_0\log_2(n+1)\) steps. If
\(\kappa=\lfloor\log_2n\rfloor\ge2\cdot10^{18}\), then
\begin{equation}\label{eq:seeds-rate}
 \frac{0.51}{v(n)}\le244\,
 \frac{\bigl(\log(4096\kappa)\bigr)^{6/11}}{\kappa^{2/11}} .
\end{equation}
Hence, for every \(n\ge2\), Algorithm~\(\mathcal U\) uses
\(n^{9/4+\delta(n)}\) steps with \(\delta(n)=O\bigl((\log\log n)^{6/11}(\log n)^{-2/11}\bigr)\).
\end{theorem}

\begin{proof}
Correctness follows from \citep[Lemma~2.4]{FinCon26}, since the maps
computed in Steps~2 and~3 are those of the relative scheme
\(\mathcal S_{1/v(n)}\), whose denominator \(\mathcal D\) is invertible in
characteristic zero. Step~1 takes \(O(\log n)\) machine steps: \(\mathsf a\)
is nondecreasing, and each value \(\mathsf a(v)\) with \(v\le\kappa\) is
computed from integers with \(O(\log\kappa)\) bits. If
\(3000\,\mathsf a(1)>\kappa\), classical multiplication adds fewer than
\(2n^3\) field operations.

Otherwise let \(v=v(n)\), and let \(A\), \(B\), \(r_*\), \(m_*\), \(p\) be
the counts and the pooled exponent of \(\mathcal S_{1/v}\). Then
\(\log_2A\le\mathsf a(v)\le\kappa/(3000v)\), so \(A\le n^{1/(3000v)}\).

\emph{Step~2.} By Proposition~\ref{prop:coefficients}(c), the generator
performs at most \(A^{46}\) bit operations, and it stores at most that
many bits. A fixed multitape Turing machine executes it, and reads its
output in Step~3, in at most \(c_{\rm TM}A^{92}\) steps, for a constant
\(c_{\rm TM}\): it keeps the
stored integers on one tape, brings the operands of each integer operation
to its work tapes by one scan of that tape, and performs schoolbook
arithmetic there in a constant times the counted bit operations. Since
\(92/3000<0.031\), this is at most \(c_{\rm TM}\,n^{0.031/v}\) steps.

\emph{Step~3.} The three maps have at most
\(3r_*A^2m_*B^2\le3A^{83/4}\) coefficients, by
Lemma~\ref{lem:resources}(b),(c). Each numerator, and \(\mathcal D\), has
at most \(\beta_\eps\le A^{1/k}\) bits
(Proposition~\ref{prop:coefficients}), so Horner's rule takes fewer than
\(2A^{1/k}\) field operations for each. With the inversion and the final
multiplications, Step~3 uses at most \(A^{22}\le n\) field operations.

\emph{Step~4.} Lemma~\ref{lem:resources} gives the hypotheses of
Lemma~\ref{lem:pooled-constants}, and \(B^{11}\le A^{22}\). By
Theorem~\ref{thm:main}, \(p<9/4+1/(2v)\). Hence Step~4 uses at most
\((304+28\log_2n)\,A^{22}n^p\le(304+28\log_2n)\,n^{9/4+(1500+22)/(3000v)}\)
field operations, and \(1522/3000<0.51\).

Since \(0.031/v\le9/4\), the four steps together use at most
\((c_0+28\log_2n)\,n^{9/4+0.51/v}\) steps, where \(c_0\) is
\(c_{\rm TM}+305\) plus the constant of Step~1.

For the rate, apply Lemma~\ref{lem:resources}(a) with \(\eps=1/v\) and
\(\Lambda(v)=\log(4096v)\): \(N_v\le2305\,v^{3/2}\Lambda(v)\) and
\(J_v\le240\,v^{3/2}\). Since \(D_v\le2\sqrt v+1\le3\sqrt v\), the number of binary
digits of \(24D_v\) is at most \(\log_2(144\sqrt v)\le\Lambda(v)\). Hence
\(\mathsf a(v)\le2\cdot2305^2\cdot240\,v^{9/2}\Lambda(v)^3
<4\cdot10^9\,v^{9/2}\Lambda(v)^3\). Put
\[
 y=\Bigl(\frac{\kappa}{1.2\cdot10^{13}\,\Lambda(\kappa)^3}\Bigr)^{2/11}.
\]
Then \(y\le\kappa^{2/11}\le\kappa\). Every integer \(v\) with
\(1\le v\le y\) satisfies \(\Lambda(v)\le\Lambda(\kappa)\), hence
\(3000\,v\,\mathsf a(v)<1.2\cdot10^{13}\,v^{11/2}\Lambda(v)^3\le\kappa\),
so \(v(n)\ge\lfloor y\rfloor\). The function \(\kappa/\Lambda(\kappa)^3\)
increases for \(\kappa\ge1\), and
\(1.2\cdot10^{13}\,\Lambda(2\cdot10^{18})^3<1.6\cdot10^{18}\). So
\(\kappa\ge2\cdot10^{18}\) gives \(y\ge1\), hence \(v(n)\ge y/2\), and
\[
 \frac{0.51}{v(n)}\le\frac{1.02}y
 =1.02\,(1.2\cdot10^{13})^{2/11}\,\frac{\Lambda(\kappa)^{6/11}}{\kappa^{2/11}}
 <244\,\frac{\Lambda(\kappa)^{6/11}}{\kappa^{2/11}},
\]
because \((1.2\cdot10^{13})^{2/11}<239\). Finally, the factor
\(c_0+28\log_2n\) adds \(O(\log\log n/\log n)\) to the exponent, which is
of smaller order.
\end{proof}

\begin{remark}[Comparison with the preceding paper]
\label{rem:uniform-comparison}
For a suitable implementation of the algorithm of
\citep[Corollary~9.7]{FinCon26}, Remark~9.8 there sketches a bound of
\(O(n^{9/4+547200/\log_2\log_2n})\) steps (its~(9.4)), an exponent below
\(3\) only when \(\log_2\log_2n>729{,}600\), and sketches why that
argument gives no better order: the seed for \(\eps=1/v\) has a
side doubly exponential in \(v\). Theorem~\ref{thm:uniform} improves the
rate to \((\log n)^{-2/11}\), up to a factor \((\log\log n)^{6/11}\). It uses that
\(\log A\) is polynomial in \(1/\eps\) and that the preprocessing costs
\(A^{O(1)}\) (Proposition~\ref{prop:coefficients}). With a preprocessing
cost \(2^{O((\log A)^c)}\) for a constant \(c>1\), the same argument would
still give a power of \(1/\log n\), with a smaller exponent. A cost
bound exponential in \(A\) would destroy the rate; the bit-length bound
of \citep[Theorem~9.3]{FinCon26} is already above \(2^A\) for the \(A\)
of that paper. The exponent \(2/11\) balances the excess
\(1/(2v)\) of the pooled exponent against \(\log(A^{22})/\log n\) with
\(\log A=O(v^{9/2}\log^3v)\); we do not know whether it can be improved.
\end{remark}

\begin{remark}[The seed route]\label{rem:seed-route}
Following the proof of \citep[Corollary~9.7]{FinCon26} instead, with the
seed of Theorem~\ref{thm:seed} and block recursion, costs at most
\(7n_0^{17/2}n^{9/4+\eps}\) operations \citep[Lemma~9.4]{FinCon26}, with
\(\log n_0=O(\eps^{-11/2}\Lambda^3)\). Balancing gives only
\(\delta(n)=O\bigl((\log\log n)^{6/13}(\log n)^{-2/13}\bigr)\), because of
the factor \(1/\eps\) between \(\log n_0\) and \(\log A\), which
Remark~\ref{rem:seed-factor} shows to be necessary for this seed.
\end{remark}
